\section*{Hoofdstuk \ref{ch:b1:predominant-reaction} --- Zuur-base-evenwichten en de methode van de predominante reactie}

\begin{solution}{exo:b1:predominant-reaction:1}
Geconjugeerde basen: \ce{HSO4-}, \ce{CO3^{2-}}, \ce{NH3}, \ce{OH-}.
Geconjugeerde zuren: \ce{NH4+}, \ce{H2PO4-}, \ce{H2O}, \ce{H3O+}.
\omterm{def:b1:predominant-reaction:ampholyte}{Amfolyten}: \ce{HCO3-}, \ce{H2O}, \ce{HPO4^{2-}} (en \ce{HSO4-}, base van
zwavelzuur en zuur van zijn tweede koppel).
\end{solution}

\begin{solution}{exo:b1:predominant-reaction:2}
\ce{HF} ($K_a = \num{6.9e-4}$) $>$ \ce{CH3COOH} (\num{1.7e-5}) $>$ \ce{HClO}
(\num{2.8e-8}) $>$ \ce{NH4+} (\num{5.6e-10}).
\end{solution}

\begin{solution}{exo:b1:predominant-reaction:3}
\omterm{def:b1:predominant-reaction:strong-acid}{Sterk zuur}: $\mathrm{pH} = 2.00$. \omterm{def:b1:predominant-reaction:strong-acid}{Sterke base}: $[\ce{OH-}] = 10^{-2}$,
$\mathrm{pH} = 14.00 - 2.00 = 12.00$.
\end{solution}

\begin{solution}{exo:b1:predominant-reaction:4}
\ce{CH3COOH} onder pH 4.76, \ce{CH3COO-} erboven. Bij pH 3 overheerst de
zure vorm (98\,\%); in bloed (pH 7.4) het ethanoaation. Bij pH 3.76 is
de fractie zuur $1/(1 + 10^{-1}) = 0.91$.
\end{solution}

\begin{solution}{exo:b1:predominant-reaction:5}
\omterm{def:b1:predominant-reaction:predominant-reaction}{Predominante reactie} \ce{NH3 + H2O <=> NH4+ + OH-}, $K = 10^{-(14.00 -
9.25)} = 10^{-4.75}$. Als er weinig \ce{NH3} reageert: $[\ce{OH-}] = \sqrt{K C} =
\qty{1.3e-3}{mol/L}$, $\mathrm{pH} = 11.13$. Controle: $[\ce{NH4+}] = 1.3\,\%$
van $C$, en $\mathrm{pH} > \mathrm{p}K_a + 1$: geldig.
\end{solution}

\begin{solution}{exo:b1:predominant-reaction:6}
$\mathrm{pH} = \frac12(4.76 + 2.00) = 3.38$ (de exacte waarde is 3.39);
$\alpha = 10^{-3.38}/0.010 = 0.042$: 4\,\% gedissocieerd, onder 10\,\%,
en $\mathrm{pH} < \mathrm{p}K_a - 1$: geldig.
\end{solution}

\begin{solution}{exo:b1:predominant-reaction:7}
\ce{HF + OH- -> F- + H2O}, $K = 10^{14.00 - 3.16} = 10^{10.84}$: aflopend.
De \omterm{def:b1:predominant-reaction:predominant-reaction}{equivalente oplossing} bevat \qty{0.050}{mol/L} van zowel \ce{HF} als
\ce{F-}, een buffer: $\mathrm{pH} = \mathrm{p}K_a = 3.16$.
\end{solution}

\begin{solution}{exo:b1:predominant-reaction:8}
$\mathrm{pH} = 4.76$. Het zoutzuur zet \qty{0.010}{mol} ethanoaat om in
zuur: $\mathrm{pH} = 4.76 + \log(0.040/0.060) = 4.58$, een verandering
van 0.18. In zuiver water zou de pH van 7.00 naar 2.00 dalen.
\end{solution}

\begin{solution}{exo:b1:predominant-reaction:9}
In water reageren beide zuren volledig tot \ce{H3O+}: het oplosmiddel
nivelleert ze. Ethaanzuur is een veel zwakkere base dan water, zodat de
twee zuren er maar gedeeltelijk mee reageren, in verschillende mate, en
hun sterkten vergeleken kunnen worden. Het sterkste zuur dat in water
kan bestaan, is \ce{H3O+}.
\end{solution}

\begin{solution}{exo:b1:predominant-reaction:10}
\ce{CO3^{2-} + H2O <=> HCO3- + OH-}, $K = 10^{-(14.00 - 10.33)} = 10^{-3.67}$;
$[\ce{OH-}] = \sqrt{10^{-3.67} \times 0.10} = \qty{4.6e-3}{mol/L}$ (4.6\,\%
van $C$), $\mathrm{pH} = 11.67$. De tweede basefunctie, \ce{HCO3- + H2O <=>
CO2 + H2O + OH-}, heeft $K = 10^{-7.63}$: ze geeft $[\ce{CO2}] \approx
10^{-7.6}$, verwaarloosbaar.
\end{solution}

\begin{solution}{exo:b1:predominant-reaction:11}
$[\mathrm{A}] = [\ce{H3O+}] = \alpha C$ en $[\mathrm{HA}] = (1 - \alpha)C$,
dus $K_a = \alpha^2C/(1 - \alpha)$. Met $C = 10^{-3}$: $\alpha^2 +
0.0174\,\alpha - 0.0174 = 0$, $\alpha = 0.12$. Het zuur is voor 12\,\%
gedissocieerd: de hypothese van een weinig gedissocieerd zuur ($\alpha <
10\,\%$) faalt, en de vierkantsvergelijking moet opgelost worden.
\end{solution}

\begin{solution}{exo:b1:predominant-reaction:12}
Het zoutzuur legt $[\ce{H3O+}] = \qty{0.010}{mol/L}$ vast: pH 2.00. Dan
is
\[
  [\ce{CH3COO-}] = \frac{K_a[\ce{CH3COOH}]}{[\ce{H3O+}]} = \frac{1.74 \times 10^{-5}
  \times 0.10}{0.010} = \qty{1.7e-4}{mol/L},
\]
verwaarloosbaar naast 0.010: het
\ce{H3O+} van het \omterm{def:b1:predominant-reaction:strong-acid}{sterke zuur} dringt het evenwicht van het \omterm{def:b1:predominant-reaction:strong-acid}{zwakke zuur}
terug.
\end{solution}

\begin{solution}{pb:b1:predominant-reaction:1}
\textbf{1.} \ce{H3PO4}/\ce{H2PO4-}, \ce{H2PO4-}/\ce{HPO4^{2-}},
\ce{HPO4^{2-}}/\ce{PO4^{3-}}; bijvoorbeeld \ce{H3PO4 <=> H2PO4- + H+}.
\textbf{2.} Grenzen bij 2.15, 7.21 en 12.34.
\textbf{3.} pH 1: \ce{H3PO4}; 5: \ce{H2PO4-}; 9: \ce{HPO4^{2-}}; 13:
\ce{PO4^{3-}}.
\textbf{4.} \ce{H2PO4-} en \ce{HPO4^{2-}}.
\textbf{5.} Elk proton laat een ion met een hogere negatieve lading
achter, dat het volgende proton steviger vasthoudt.
\textbf{6.} \ce{H3PO4 + H2O <=> H2PO4- + H3O+}; de formule
$\frac12(\mathrm{p}K_a + \mathrm{p}C) = 1.58$ schendt $\mathrm{pH} <
\mathrm{p}K_a - 1 = 1.15$: ongeveer een vijfde van het zuur reageert.
\textbf{7.} $x^2 + K_{a1}x - K_{a1}C = 0$ met $K_{a1} = 10^{-2.15}$: $x =
\qty{0.0233}{mol/L}$, $\mathrm{pH} = 1.63$.
\textbf{8.} \ce{2H2PO4- <=> H3PO4 + HPO4^{2-}}, $K = 10^{2.15 - 7.21} =
10^{-5.06}$.
\textbf{9.} $\mathrm{pH} = \frac12(2.15 + 7.21) = 4.68$.
\textbf{10.} $\mathrm{pH} = \frac12(7.21 + 12.34) = 9.78$.
\textbf{11.} Geen enkele: 1.63, 4.68 en 9.78; er is een mengsel van twee
ervan nodig.
\textbf{12.} \ce{H3PO4 + OH- -> H2PO4- + H2O}, $K = 10^{14.00 - 2.15} =
10^{11.85}$: \omterm{def:b1:extent-q-and-k:quantitative}{kwantitatief}.
\textbf{13.} \qty{0.100}{mol/L} \ce{H2PO4-}: pH 4.68.
\textbf{14.} De volgende \qty{0.050}{mol} hydroxide zet er de helft van
om: \qty{0.050}{mol/L} van zowel \ce{H2PO4-} als \ce{HPO4^{2-}}, pH 7.21.
\textbf{15.} \qty{0.100}{mol/L} \ce{HPO4^{2-}}: pH 9.78.
\textbf{16.} De pH stijgt vanaf 1.63, maakt een sprong rond $n = 0.100$
(naar 4.68), stijgt langzaam door 7.21 bij $n = 0.150$, maakt een sprong
rond $n = 0.200$ (naar 9.78), passeert 12.34 bij ongeveer $n = 0.250$ en
bereikt ongeveer 12.6 bij $n = 0.300$ (\qty{0.100}{mol/L} \ce{PO4^3-}:
$x^2/(0.100 - x) = 10^{-1.66}$, $x = 0.037$, pH 12.57).
\textbf{17.} Ze bevat gelijke hoeveelheden van het zuur en de base van een
koppel met $\mathrm{p}K_a$ 7.21 (\cref{prop:b1:predominant-reaction:buffer}).
\textbf{18.} $10^{7.20 - 7.21} = 0.977$.
\textbf{19.} $n(\ce{H2PO4-}) = 0.100 - x$, $n(\ce{HPO4^{2-}}) = x$.
\textbf{20.} $x/(0.100 - x) = 0.977$, $x = \qty{0.0494}{mol}$.
\textbf{21.} Het zuur zet \qty{0.0010}{mol} \ce{HPO4^{2-}} om:
$\mathrm{pH} = 7.21 + \log(0.0484/0.0516) = 7.18$, een verandering van
$-0.02$. In zuiver water: van 7.00 naar 3.00.
\textbf{22.} Nee: de pH hangt af van de verhouding van de hoeveelheden,
die door verdunning niet verandert (zolang de \omterm{def:b1:extent-q-and-k:activity}{activiteiten} dicht bij de
concentraties blijven).
\textbf{23.} $0.100 + 0.0494 = \qty{0.149}{mol}$ hydroxide:
\textbf{\qty{149}{mL}} van de oplossing van \qty{1.00}{mol/L}, aangevuld
tot \qty{1.00}{L}.
\end{solution}
